\renewcommand{\thesection}{\arabic{section}}\renewcommand{\theHsection}{chapter.\arabic{section}}\setcounter{section}{3}\section{一元微分学的顶峰——Taylor定理}
\noindent\begin{minipage}{\textwidth}\ProblemHeading{例题}{4.1（改写）}
圆的半径由 $r>0$ 变为 $r+h>0$。求面积的改变量及其关于 $h$ 的线性近似，并计算近似的误差。当 $h>0$ 时，说明它们与圆环面积、长为 $2\pi r$ 且宽为 $h$ 的矩形面积的关系。
\end{minipage}\par\vfill
\noindent\begin{minipage}{\textwidth}\ProblemHeading{例题}{4.2}
用微分作近似计算：$\sin(30^\circ5^{\prime})$。
\end{minipage}\par\vfill
\noindent\begin{minipage}{\textwidth}\ProblemHeading{例题}{4.3}
求 $\mathrm d(x^5)$。
\end{minipage}\par\vfill
\newpage
\noindent\begin{minipage}{\textwidth}\ProblemHeading{例题}{4.4}
求 $\mathrm d\bigl(\mathrm e^{x^2-1/x}\bigr)$。
\end{minipage}\par\vfill
\noindent\begin{minipage}{\textwidth}\ProblemHeading{习题}{4.5}
求下列函数 $y$ 关于 $x$ 的微分；出现的参数为常数，在自然定义域内讨论：
\SubQ{1}{$y=1/x$。}
\SubQ{2}{$y=\arctan(ax+b)$。}
\SubQ{3}{$y=\sin x-x\cos x$。}
\end{minipage}\par\vfill
\noindent\begin{minipage}{\textwidth}\ProblemHeading{习题}{4.6}
求下列函数 $y$ 关于 $x$ 的微分；出现的参数为常数，在自然定义域内讨论：
\SubQ{1}{$y=\ln(x+\sqrt{x^2+a^2})$。}
\end{minipage}\par\vfill
\newpage
\noindent\begin{minipage}{\textwidth}\ProblemHeading{习题}{4.7}
在括号中填入一个函数，使等式在适当的定义区间内成立；$a$ 为常数：
\SubQ{1}{$\mathrm d(\qquad)=\mathrm dx/x$。}
\SubQ{2}{$\mathrm d(\qquad)=\mathrm dx/(1+x^2)$。}
\SubQ{3}{$\mathrm d(\qquad)=\mathrm dx/\sqrt{x}$。}
\end{minipage}\par\vfill
\noindent\begin{minipage}{\textwidth}\ProblemHeading{习题}{4.8}
在括号中填入一个函数，使等式在适当的定义区间内成立；$a$ 为常数：
\SubQ{1}{$\mathrm d(\qquad)=(2x+1)\,\mathrm dx$。}
\SubQ{2}{$\mathrm d(\qquad)=(\cos x+\sin x)\,\mathrm dx$。}
\SubQ{3}{$\mathrm d(\qquad)=\mathrm dx/\sqrt{1+x^2}$。}
\end{minipage}\par\vfill
\noindent\begin{minipage}{\textwidth}\ProblemHeading{习题}{4.9}
在括号中填入一个函数，使等式在适当的定义区间内成立；$a$ 为常数：
\SubQ{1}{$\mathrm d(\qquad)=\mathrm e^{-ax}\,\mathrm dx$。}
\SubQ{2}{$\mathrm d(\qquad)=\cos x\sin x\,\mathrm dx$。}
\SubQ{3}{$\mathrm d(\qquad)=\mathrm dx/(x\ln x)$。}
\end{minipage}\par\vfill
\newpage
\noindent\begin{minipage}{\textwidth}\ProblemHeading{习题}{4.10}
在括号中填入一个函数，使等式在适当的定义区间内成立；$a$ 为常数：
\SubQ{1}{$\mathrm d(\qquad)=x\,\mathrm dx/\sqrt{x^2+a^2}$。}
\SubQ{2}{$\mathrm d(\qquad)=\cos^2x\sin x\,\mathrm dx$。}
\SubQ{3}{$\mathrm d(\qquad)=\sin^2x\,\mathrm dx$。}
\end{minipage}\par\vfill
\noindent\begin{minipage}{\textwidth}\ProblemHeading{习题}{4.11}
设 $u,v,w$ 均为 $x$ 的可微函数。求下列 $y$ 关于 $x$ 的微分，在函数有定义的范围内讨论：
\SubQ{1}{$y=uvw$。}
\SubQ{2}{$y=u/v^2$。}
\SubQ{3}{$y=1/\sqrt{u^2+v^2+w^2}$。}
\end{minipage}\par\vfill
\noindent\begin{minipage}{\textwidth}\ProblemHeading{习题}{4.12}
设 $u,v,w$ 均为 $x$ 的可微函数。求下列 $y$ 关于 $x$ 的微分，在函数有定义的范围内讨论：
\SubQ{1}{$y=\arctan(u/(vw))$。}
\SubQ{2}{$y=\ln\sqrt{u^2+v^2+w^2}$。}
\end{minipage}\par\vfill
\newpage
\noindent\begin{minipage}{\textwidth}\ProblemHeading{习题}{4.13}
对下列方程作微分，并在能局部确定可导函数 $y=y(x)$ 的点求出 $y^{\prime}$；参数均为常数：
\SubQ{1}{$x=y+\mathrm e^y$。}
\SubQ{2}{$x=y+\ln y$。}
\SubQ{3}{$x^2/a^2+y^2/b^2=1$，$a,b\ne0$。}
\end{minipage}\par\vfill
\noindent\begin{minipage}{\textwidth}\ProblemHeading{习题}{4.14}
对下列方程作微分，并在能局部确定可导函数 $y=y(x)$ 的点求出 $y^{\prime}$；参数均为常数：
\SubQ{1}{$\sqrt{x}+\sqrt{y}=\sqrt a$，$a>0$。}
\SubQ{2}{$y^2/x=\sqrt{x^2+y^2}$。}
\end{minipage}\par\vfill
\noindent\begin{minipage}{\textwidth}\ProblemHeading{习题}{4.15}
利用微分作近似计算：
\SubQ{1}{$\sqrt[4]{80}$。}
\SubQ{2}{$\sin29^\circ$。}
\SubQ{3}{$\arctan1.05$。}
\end{minipage}\par\vfill
\newpage
\noindent\begin{minipage}{\textwidth}\ProblemHeading{习题}{4.16}
利用微分作近似计算：
\SubQ{1}{$\lg11$（常用对数）。}
\end{minipage}\par\vfill
\noindent\begin{minipage}{\textwidth}\ProblemHeading{例题}{4.17}
求 $f(x)=\mathrm e^x$ 在 $0$ 处带 Peano 余项的 Maclaurin 展开式，写至 $n$ 次，$n\in\mathbb N^*$。
\end{minipage}\par\vfill
\noindent\begin{minipage}{\textwidth}\ProblemHeading{例题}{4.18}
求 $f(x)=\sin x$ 在 $0$ 处带 Peano 余项的 Maclaurin 展开式，写至 $2n$ 次，$n\in\mathbb N^*$。
\end{minipage}\par\vfill
\newpage
\noindent\begin{minipage}{\textwidth}\ProblemHeading{例题}{4.19（改写）}
证明当 $x\to0$ 时，$\cos x=\sum_{k=0}^n(-1)^k x^{2k}/(2k)!+o(x^{2n+1})$，其中 $n\in\mathbb N^*$。
\end{minipage}\par\vfill
\noindent\begin{minipage}{\textwidth}\ProblemHeading{例题}{4.20}
求 $f(x)=\ln(1+x)$ 在 $0$ 处带 Peano 余项的 Maclaurin 展开式，写至 $n$ 次，$n\in\mathbb N^*$。
\end{minipage}\par\vfill
\noindent\begin{minipage}{\textwidth}\ProblemHeading{例题}{4.21}
设 $\lambda\in\mathbb R$。求 $f(x)=(1+x)^\lambda$（$x>-1$）在 $0$ 处带 Peano 余项的 Maclaurin 展开式，写至 $n$ 次，$n\in\mathbb N^*$。
\end{minipage}\par\vfill
\newpage
\noindent\begin{minipage}{\textwidth}\ProblemHeading{例题}{4.22}
求 $f(x)=\arctan x$ 在 $0$ 处带 Peano 余项的 Maclaurin 展开式，写至 $2n+2$ 次，$n\in\mathbb N^*$。
\end{minipage}\par\vfill
\noindent\begin{minipage}{\textwidth}\ProblemHeading{习题}{4.23}
求 $\arcsin x$ 在 $0$ 处带 Peano 余项的 Maclaurin 展开式，并给出一般系数。
\end{minipage}\par\vfill
\noindent\begin{minipage}{\textwidth}\ProblemHeading{例题}{4.24}
计算 $\lim_{x\to0}\dfrac{\mathrm e^x\sin x-x(1+x)}{\sin^3x}$。
\end{minipage}\par\vfill
\newpage
\noindent\begin{minipage}{\textwidth}\ProblemHeading{例题}{4.25}
计算 $\lim_{x\to\infty}\left(x-x^2\ln(1+1/x)\right)$。
\end{minipage}\par\vfill
\noindent\begin{minipage}{\textwidth}\ProblemHeading{例题}{4.26}
求 $\lim_{x\to0}\dfrac{\sin x-\arctan x}{\tan x-\sin x}$。
\end{minipage}\par\vfill
\noindent\begin{minipage}{\textwidth}\ProblemHeading{例题}{4.27}
设 $a\in\mathbb R$。求 $\lim_{n\to\infty}\prod_{k=1}^n\cos\dfrac{ka}{n\sqrt n}$。
\end{minipage}\par\vfill
\newpage
\noindent\begin{minipage}{\textwidth}\ProblemHeading{习题}{4.28}
计算下列极限（无符号的 $x\to\infty$ 按 $|x|\to+\infty$ 理解，若双侧极限不存在须说明）：
\SubQ{1}{$\lim_{x\to0}\dfrac{\cos x-\mathrm e^{-x^2/2}}{x^4}$。}
\SubQ{2}{$\lim_{x\to0}\dfrac{\mathrm e^{\alpha x}-\mathrm e^{\beta x}}{\sin(\alpha x)-\sin(\beta x)}$，$\alpha\ne\beta$。}
\SubQ{3}{$\lim_{x\to\infty}\left((x^3-x^2+x/2)\mathrm e^{1/x}-\sqrt{x^6+1}\right)$。}
\end{minipage}\par\vfill
\noindent\begin{minipage}{\textwidth}\ProblemHeading{习题}{4.29}
计算下列极限（无符号的 $x\to\infty$ 按 $|x|\to+\infty$ 理解，若双侧极限不存在须说明）：
\SubQ{1}{$\lim_{x\to0}\dfrac{a^x+a^{-x}-2}{x^2}$，$a>0$。}
\end{minipage}\par\vfill
\noindent\begin{minipage}{\textwidth}\ProblemHeading{习题}{4.30}
设函数 $f$ 在 $x_0$ 附近可以表示为 $f(x)=\sum_{k=0}^n a_k(x-x_0)^k+o((x-x_0)^n)$。是否必定有 $a_k=f^{(k)}(x_0)/k!$（$k=0,1,\ldots,n$）？说明理由。
\end{minipage}\par\vfill
\newpage
\noindent\begin{minipage}{\textwidth}\ProblemHeading{习题}{4.31}
设 $\alpha\ge-1$。计算 $\lim_{n\to\infty}\left[\prod_{k=1}^n(1+k/n^{\alpha+2})\right]^{n^\alpha}$。
\end{minipage}\par\vfill
\noindent\begin{minipage}{\textwidth}\ProblemHeading{习题}{4.32}
令 $f(x)=\lim_{n\to\infty}n^x\left((1+1/(n+1))^{n+1}-(1+1/n)^n\right)$，以极限存在且有限的实数 $x$ 为定义域。求 $f$ 的定义域和值域。
\end{minipage}\par\vfill
\noindent\begin{minipage}{\textwidth}\ProblemHeading{习题}{4.33}
证明 $\lim_{n\to\infty}\left(\dfrac{n^4}{\mathrm e}(1+1/n)^n-n^4+\dfrac{n^3}{2}-\dfrac{11n^2}{24}+\dfrac{7n}{16}\right)=\dfrac{2447}{5760}$。
\end{minipage}\par\vfill
\newpage
\noindent\begin{minipage}{\textwidth}\ProblemHeading{习题}{4.34}
证明多项式 $\sum_{k=1}^n\dfrac{(2x-x^2)^k-2x^k}{k}$ 能被 $x^{n+1}$ 整除，其中 $n\in\mathbb N^*$。
\end{minipage}\par\vfill
\noindent\begin{minipage}{\textwidth}\ProblemHeading{习题}{4.35（校订）}
设 $x_1=\sin x_0>0$，$x_{n+1}=\sin x_n$（$n\ge1$）。证明 $\lim_{n\to\infty}\sqrt{n/3}\,x_n=1$。
\end{minipage}\par\vfill
\noindent\begin{minipage}{\textwidth}\ProblemHeading{习题}{4.36}
设 $y_1=c>0$，$y_{n+1}/(n+1)=\ln(1+y_n/n)$（$n\ge1$）。证明 $\lim_{n\to\infty}y_n=2$。
\end{minipage}\par\vfill
\newpage
\noindent\begin{minipage}{\textwidth}\ProblemHeading{例题}{4.37（改写）}
写出 $\mathrm e^x$ 在 $0$ 处 $n$ 次 Maclaurin 展开式的 Lagrange 余项，并将其中的中间点写成 $\theta x$（$0<\theta<1$），$n\in\mathbb N^*$。
\end{minipage}\par\vfill
\noindent\begin{minipage}{\textwidth}\ProblemHeading{例题}{4.38（改写）}
将 $\sin x$ 在 $0$ 处展开至 $2n$ 次，求其 Lagrange 余项 $R_{2n}(x)$，并用中间点 $\theta x$（$0<\theta<1$）表示，$n\in\mathbb N^*$。
\end{minipage}\par\vfill
\noindent\begin{minipage}{\textwidth}\ProblemHeading{例题}{4.39（改写）}
将 $\cos x$ 在 $0$ 处展开至 $2n+1$ 次，求其 Lagrange 余项 $R_{2n+1}(x)$，并用中间点 $\theta x$（$0<\theta<1$）表示，$n\in\mathbb N^*$。
\end{minipage}\par\vfill
\newpage
\noindent\begin{minipage}{\textwidth}\ProblemHeading{例题}{4.40（改写）}
将 $\ln(1+x)$（$x>-1$）在 $0$ 处展开至 $n$ 次，求其 Lagrange 余项，并用中间点 $\theta x$（$0<\theta<1$）表示，$n\in\mathbb N^*$。
\end{minipage}\par\vfill
\noindent\begin{minipage}{\textwidth}\ProblemHeading{例题}{4.41（改写）}
设 $\lambda\in\mathbb R$，$n\in\mathbb N^*$。写出 $(1+x)^\lambda$（$x>-1$）在 $0$ 处至 $n$ 次的 Maclaurin 展开式，分别求其 Lagrange 余项和 Cauchy 余项；再利用 Cauchy 余项估计 $|x|<1$ 时的展开误差。
\end{minipage}\par\vfill
\noindent\begin{minipage}{\textwidth}\ProblemHeading{例题}{4.42}
利用 Taylor 展开求 $2^{1/5}$ 的近似值，并估计误差。
\end{minipage}\par\vfill
\newpage
\noindent\begin{minipage}{\textwidth}\ProblemHeading{例题}{4.43}
在 $[0,\pi]$ 上，用九次多项式 $x-x^3/3!+x^5/5!-x^7/7!+x^9/9!$ 逼近 $\sin x$。给出一个对整个区间成立的误差上界。
\end{minipage}\par\vfill
\noindent\begin{minipage}{\textwidth}\ProblemHeading{例题}{4.44}
设 $f$ 在 $[0,1]$ 上有二阶导数，且 $|f(0)|\le1$、$|f(1)|\le1$、$|f^{\prime\prime}(x)|\le2$。证明对所有 $x\in[0,1]$ 均有 $|f^{\prime}(x)|\le3$。
\end{minipage}\par\vfill
\noindent\begin{minipage}{\textwidth}\ProblemHeading{例题}{4.45}
设 $f$ 在 $[a,+\infty)$ 上有三阶导数，且 $\lim_{x\to+\infty}f(x)$ 与 $\lim_{x\to+\infty}f^{\prime\prime\prime}(x)$ 均存在且有限。证明 $\lim_{x\to+\infty}f^{\prime}(x)=\lim_{x\to+\infty}f^{\prime\prime}(x)=\lim_{x\to+\infty}f^{\prime\prime\prime}(x)=0$。
\end{minipage}\par\vfill
\newpage
\noindent\begin{minipage}{\textwidth}\ProblemHeading{习题}{4.46}
将多项式 $1+2x+3x^2+4x^3+5x^4$ 按 $x+1$ 的幂展开。
\end{minipage}\par\vfill
\noindent\begin{minipage}{\textwidth}\ProblemHeading{习题}{4.47}
按指定次数写出函数 $f$ 在 $x=0$ 处的 Taylor 多项式：
\SubQ{1}{$f(x)=(1+x+x^2)/(1-x+x^2)$，写至四次。}
\SubQ{2}{$f(x)=\mathrm e^{2x-x^2}$，写至五次。}
\SubQ{3}{$f(x)=\ln\cos x$，写至六次。}
\end{minipage}\par\vfill
\noindent\begin{minipage}{\textwidth}\ProblemHeading{习题}{4.48}
按指定次数写出函数 $f$ 在 $x=0$ 处的 Taylor 多项式：
\SubQ{1}{$f(x)=\tan x$，写至五次。}
\SubQ{2}{$f(x)=1/\sqrt{1-x^2}$，写至六次。}
\end{minipage}\par\vfill
\newpage
\noindent\begin{minipage}{\textwidth}\ProblemHeading{习题}{4.49}
按指定次数写出函数 $f$ 在指定点 $x_0$ 处的 Taylor 多项式，$n\in\mathbb N^*$：
\SubQ{1}{$f(x)=\sin x$，$x_0=\pi/2$，写至 $2n$ 次。}
\SubQ{2}{$f(x)=\cos x$，$x_0=\pi$，写至 $2n$ 次。}
\SubQ{3}{$f(x)=\mathrm e^x$，$x_0=1$，写至 $n$ 次。}
\end{minipage}\par\vfill
\noindent\begin{minipage}{\textwidth}\ProblemHeading{习题}{4.50}
按指定次数写出函数 $f$ 在指定点 $x_0$ 处的 Taylor 多项式，$n\in\mathbb N^*$：
\SubQ{1}{$f(x)=\ln x$，$x_0=2$，写至 $n$ 次。}
\SubQ{2}{$f(x)=x/(1+x^2)$，$x_0=0$，写至 $2n+1$ 次。}
\end{minipage}\par\vfill
\noindent\begin{minipage}{\textwidth}\ProblemHeading{习题}{4.51}
设 $f,g$ 在 $(-1,1)$ 上无限次可导，且对 $|x|<1$、$n=0,1,2,\ldots$ 有 $|f^{(n)}(x)-g^{(n)}(x)|\le n!|x|$。证明 $f=g$。
\end{minipage}\par\vfill
\newpage
\noindent\begin{minipage}{\textwidth}\ProblemHeading{习题}{4.52}
证明对 $x>0$、$n\in\mathbb N^*$ 有 $\sum_{k=1}^{2n}(-1)^{k-1}x^k/k<\ln(1+x)<\sum_{k=1}^{2n-1}(-1)^{k-1}x^k/k$。
\end{minipage}\par\vfill
\noindent\begin{minipage}{\textwidth}\ProblemHeading{习题}{4.53}
设 $n\in\mathbb N^*$，$f$ 在 $x_0$ 处有 $n+1$ 阶导数，且 $f^{(n+1)}(x_0)\ne0$。在充分小的非零 $h$ 下，将 Taylor 公式写为 \[f(x_0+h)=\sum_{k=0}^{n-1}\frac{f^{(k)}(x_0)}{k!}h^k+\frac{h^n}{n!}f^{(n)}(x_0+\theta_n h),\qquad 0<\theta_n<1.\] 证明 $\lim_{h\to0}\theta_n=1/(n+1)$。
\end{minipage}\par\vfill
\noindent\begin{minipage}{\textwidth}\ProblemHeading{习题}{4.54}
设 $a<b$，$f$ 在 $[a,b]$ 上有二阶导数，$f^{\prime}(a)=f^{\prime}(b)=0$。证明存在 $c\in(a,b)$ 使 $|f^{\prime\prime}(c)|\ge4|f(b)-f(a)|/(b-a)^2$。
\end{minipage}\par\vfill
\newpage
\noindent\begin{minipage}{\textwidth}\ProblemHeading{习题}{4.55}
设 $f$ 在 $[0,1]$ 上有二阶导数，$f(0)=f(1)=0$，且 $f$ 在此区间上的最小值为 $-1$。证明存在 $\xi\in(0,1)$ 使 $f^{\prime\prime}(\xi)\ge8$。
\end{minipage}\par\vfill
\noindent\begin{minipage}{\textwidth}\ProblemHeading{习题}{4.56（校订）}
设 $n\ge2$ 为整数，$f$ 在 $(x_0-\delta,x_0+\delta)$ 上有 $n$ 阶导数，且 $f^{(k)}(x_0)=0$（$2\le k\le n-1$），$f^{(n)}(x_0)\ne0$，$f^{(n)}$ 在 $x_0$ 处连续。当 $0<|h|<\delta$ 时，设 $f(x_0+h)-f(x_0)=hf^{\prime}(x_0+\theta h)$，其中 $0<\theta<1$。证明 $\lim_{h\to0}\theta=n^{-1/(n-1)}$。
\end{minipage}\par\vfill
\noindent\begin{minipage}{\textwidth}\ProblemHeading{习题}{4.57}
令 $P_n(x)=\sum_{k=0}^n x^k/k!$，$n\in\mathbb N^*$。证明：
\SubQ{1}{当 $n$ 为偶数时，$P_n(x)>0$ 对所有实数 $x$ 成立。}
\SubQ{2}{当 $n$ 为奇数时，$P_n$ 有唯一实零点。}
\SubQ{3}{记 $P_{2n+1}$ 的实零点为 $x_n$（$n=0,1,2,\ldots$），则数列 $\{x_n\}$ 严格递减且趋于 $-\infty$。}
\end{minipage}\par\vfill
\newpage
\noindent\begin{minipage}{\textwidth}\ProblemHeading{习题}{4.58}
设 $f$ 在 $[0,2]$ 上有二阶导数，$|f(0)|\le1$、$|f(2)|\le1$，且 $|f^{\prime\prime}|\le1$。证明在此区间上 $|f^{\prime}|\le2$，并证明常数 $2$ 不能减小。
\end{minipage}\par\vfill
\noindent\begin{minipage}{\textwidth}\ProblemHeading{习题}{4.59}
设 $f$ 在 $\mathbb R$ 上二次可导，$M_k=\sup_{x\in\mathbb R}|f^{(k)}(x)|<+\infty$（$k=0,1,2$）。证明 $M_1^2\le2M_0M_2$。
\end{minipage}\par\vfill
